% !TEX root = ../main.tex
本迪布焦型埃博拉疫情自 2026 年 5 月起在刚果（金）8 个省、66 个卫生区扩散，累计确诊约 10\,060 例、死亡约 4\,887 人，肯尼亚与尼日利亚已出现输入病例，题目要求我们从传播机制一直做到面向卫生部与国际捐助方的执行摘要。

任务一，我们建立含人际、院内与葬礼三条传播通道的双索引八舱连续时间模型（$S\to V\to E\to I\to H\to T\to F\to R$），以省与家庭双索引刻画省际流动，并用时变传播强度 $k(t)$ 描述干预随时间加强。

任务二，以 27 周全国序列、16 节点对数样条加二阶差分惩罚估计 $k(t)$ 与检出比例：样本内 $R^2=0.955$、RMSE $=34.6$ 例/周、MAPE $=6.9\%$，前 20 周拟合、后 7 周回测 RMSE $=119.5$、MAPE $=20.2\%$；得 $\theta=0.3382$、$\delta=0.486$、$\rho=0.392$，后期有效再生数为 0.893。

任务三，在统一情景口径下从第 27 周外推 26 周，基线新增 5\,672 例、报告口径死亡 2\,937 例。单措施中安全安葬最有效（2\,172 例，$-61.7\%$，30 周内收尾），接种 5\%/周为 2\,591 例（$-54.3\%$），院内感染控制 $-17.1\%$，治疗 50\% 为 4\,199 例（$-26.0\%$），而单独加快隔离与单独提高追踪率反而使新增增加 2.0\% 与 3.4\%。组合 C2 为 1\,834 例（$-67.7\%$）、C3 为 1\,590 例（$-72.0\%$）。若 7 月中旬即达到当前干预强度，同期末为 6\,884 例而非 17\,188 例，少 10\,303 例（$59.9\%$）与 4\,833 例死亡，收尾提前约 30 周。随访率由 78\% 提高到 90\%/95\%，26 周新增降至 5\,520/4\,720 例，需增约 3\,200/4\,600 名社区工作者与 6\,150/8\,700 万美元。

任务四，省际流动矩阵由重力先验加单参数标定；出行限制只改变空间分布与响应时间，不能单独压低全国总量。跨境不直接外推当地数值，而按 $R_{\mathrm{local}}\in\{0.8,1.0,1.2,1.5\}$ 判定：$R_{\mathrm{local}}\le 1$ 时输入链必然熄灭，1.2 时为 19\%--23\%，1.5 时为 3\%--4\%。

任务五，在传播强度 $\pm30\%$、检出比例 $\rho\in[0.30,0.50]$ 与干预时机联合扰动下，安葬优先于接种的排序不变，但有效再生数的结构性包络 $[0.62,1.22]$ 跨越 1，故不把“已确定低于 1”当作决策前提。

任务六，我们建议优先保障安全安葬和医院感染控制，并加强接触者追踪。对于超过 4 亿美元的资金缺口，国际援助不仅关系到当地救治，也关系到周边国家能否减少输入病例及后续防控支出。  
  
\noindent\rule{\linewidth}{0.4pt}
                                          
\noindent\textbf{关键词：}埃博拉病毒病；双索引SEIHFR 模型；时变参数；省际流动；下一代矩阵；干预组合优化；不确定性分析